% ============================================================
% 5_E_frontend_interact.tex — 前端交互与UI模块概要设计
% 编写人：E  日期：2026-07-15
% 职责：管理员端页面 / 教师端页面 / UI设计规范
% ============================================================

\section{管理员端页面概要设计}

\subsection{页面布局}

管理员端采用侧边导航+顶部栏+内容区域的经典后台管理布局。侧边导航收起时仅显示图标，展开时显示图标+文字标签。顶部栏展示系统Logo、当前管理员用户名、退出登录按钮。内容区域根据路由切换展示不同的功能页面。

管理员端包含以下核心页面：
\begin{itemize}
    \item \textbf{仪表盘页（/admin/dashboard）}：系统数据概览看板，展示用户数、论文数、模板数等核心指标；
    \item \textbf{学院管理页（/admin/colleges）}：学院信息的CRUD操作界面；
    \item \textbf{用户管理页（/admin/users）}：用户账号的查询、筛选、禁用/启用操作界面；
    \item \textbf{模板管理页（/admin/templates）}：模板列表展示、创建、编辑、启用/停用操作界面；
    \item \textbf{模板版本编辑页（/admin/templates/:id/versions）}：JSON配置编辑器，管理模板的封面字段、章节结构和格式参数。
\end{itemize}

\subsection{仪表盘页设计}

仪表盘页采用卡片网格布局，上方为四个统计指标卡片（总用户数、在编论文数、模板总数、今日导出数），下方为图表区域（论文状态分布饼图、每天导出数量趋势折线图）。所有图表数据通过独立API获取，支持按时间范围筛选。

\subsection{学院管理页设计}

学院管理页采用表格视图，上方为操作栏（新增学院按钮、搜索输入框），主体为el-table组件展示学院列表，字段包含学院名称、学院代码、关联模板数、关联用户数、创建时间。表格行末尾提供编辑和删除操作按钮。新增/编辑操作通过el-dialog弹窗实现。

\subsection{模板管理页设计}

模板管理页是管理员端配置最复杂的页面，采用左侧列表+右侧编辑的双栏布局：
\begin{itemize}
    \item 左侧模板列表以卡片形式展示模板名称、类型标签（毕业/课程/项目）、当前版本号、启用状态开关（el-switch），支持按类型和学院筛选；
    \item 右侧编辑区域在未选中模板时展示"请选择模板"占位提示，选中模板后切换为编辑界面；
    \item 编辑界面采用Tabs组件组织三个配置面板：封面字段配置（拖拽排序的字段列表）、章节结构配置（树形编辑器）、格式参数配置（表单）。
\end{itemize}

\subsection{模板版本编辑页设计}

版本编辑页右侧使用JSON编辑器或可视化表单进行配置编辑。封面字段配置区支持拖拽添加/删除/排序封面元素，每个元素可配置字段标签、是否必填、默认值。章节结构采用el-tree组件，支持拖拽调整层级、右键菜单新增/删除节点。格式参数采用el-form配合el-color-picker（颜色选择）、el-slider（行距/边距调整）等丰富控件，参数调整后右侧实时预览展示渲染效果。

\section{教师端页面概要设计}

\subsection{页面布局}

教师端同样采用侧边导航+顶部栏+内容区域布局。侧边导航包含批阅论文列表和批阅历史两个主要入口。教师端整体风格偏阅读和审阅，内容区域以白底卡片为主，减少视觉干扰。

教师端包含以下核心页面：
\begin{itemize}
    \item \textbf{批阅任务列表页（/teacher/reviews）}：展示待批阅和已批阅的论文列表；
    \item \textbf{论文批阅页（/teacher/reviews/:id）}：论文全文阅读+批注添加+评分的核心操作页面；
    \item \textbf{批阅历史页（/teacher/history）}：展示教师历次批阅记录汇总。
\end{itemize}

\subsection{批阅任务列表页设计}

批阅任务列表采用表格视图，展示字段包括论文标题、学生姓名、所属学院、提交时间、当前状态（待批阅/已批阅/已退回）。表格支持按学院筛选和按提交时间排序。状态列使用不同颜色的el-tag展示：待批阅为橙色、已批阅为绿色、已退回为灰色。新增提交的论文在列表顶部显示，并用闪烁的"新"标签提示教师。

\subsection{论文批阅页设计}

论文批阅页采用三栏布局，是本系统交互最复杂的页面之一：
\begin{enumerate}
    \item \textbf{左侧栏（章节导航，宽度240px）}：使用el-tree展示论文的完整章节树，当前批阅章节高亮显示。教师点击章节名称切换预览内容。章节标题旁标注该章节的批注数量；
    \item \textbf{中间栏（论文预览，自适应宽度）}：以HTML渲染方式展示论文章节内容，模拟A4纸张样式。页面支持滚动浏览。教师拖拽选中文字后，选中区域以主题色高亮显示，并弹出浮动工具栏（添加批注按钮）。已添加批注的文字以黄色下划线标记，悬停时显示批注预览；
    \item \textbf{右侧栏（批注面板，宽度320px）}：分为上下两部分。上部为批注列表区，展示当前章节的全部批注（按创建时间排序），每条批注显示批注内容摘要和创建时间。下部为批注编辑区，当教师选中文字后自动填充selected\_text预览，并提供textarea输入批注内容。
\end{enumerate}

\subsection{评分对话框设计}

教师完成批注后，点击页面顶部的"评阅完成"按钮触发评分对话框。对话框采用el-dialog组件，宽度600px，包含以下控件：
\begin{itemize}
    \item 评语编辑器（el-input type="textarea"，支持富文本格式，最大5000字符）；
    \item 评分滑块（el-slider，范围0-100，步进1，显示当前分值大号数字）；
    \item 等级展示（根据分数自动换算并显示，如"85分→良"，带颜色标签）；
    \item 操作类型单选组（el-radio-group）：通过（REVIEWED）/ 退回（RETURNED）；
    \item 退回原因输入框（条件显示，仅当选择"退回"时出现，el-input，最小10字符）；
    \item 确认提交按钮和取消按钮。
\end{itemize}

\section{UI设计规范}

\subsection{颜色系统}

为保证系统视觉一致性，定义以下颜色令牌：

\begin{table}[htbp]
\centering
\caption{UI颜色系统规范}
\label{tab:ui_color}
\begin{tabular}{|c|c|c|c|}
\hline
\textbf{用途} & \textbf{色值} & \textbf{Element Plus变量} & \textbf{适用范围} \\
\hline
主色 & \#409EFF & --el-color-primary & 按钮、链接、导航选中态 \\
\hline
成功 & \#67C23A & --el-color-success & 成功提示、已完成标签 \\
\hline
警告 & \#E6A23C & --el-color-warning & 警告提示、待处理标签 \\
\hline
危险 & \#F56C6C & --el-color-danger & 错误提示、删除操作、已退回标签 \\
\hline
信息 & \#909399 & --el-color-info & 辅助文字、禁用状态 \\
\hline
背景 & \#F5F7FA & — & 页面背景色、卡片背景色 \\
\hline
\end{tabular}
\end{table}

\subsection{字号层级}

\begin{table}[htbp]
\centering
\caption{字号层级规范}
\label{tab:font_size}
\begin{tabular}{|c|c|c|}
\hline
\textbf{层级} & \textbf{字号} & \textbf{应用场景} \\
\hline
页面标题 & 20px / 18pt & 各功能页面的主标题 \\
\hline
卡片标题 & 16px / 14pt & 卡片组件标题、Dialog标题 \\
\hline
正文/表头 & 14px / 12pt & 表格内容、表单标签、段落文字 \\
\hline
辅助文字 & 12px / 10.5pt & 提示信息、时间戳、标签文字 \\
\hline
\end{tabular}
\end{table}

\subsection{间距规范}

\begin{itemize}
    \item \textbf{页面内边距}：内容区域与页面边缘保持24px间距；
    \item \textbf{卡片内边距}：卡片容器内边距20px；
    \item \textbf{组件间距}：相邻表单控件之间间距16px；
    \item \textbf{标签与控件间距}：表单标签与输入框之间间距12px；
    \item \textbf{段落间距}：卡片内多段文字间距8px。
\end{itemize}

\subsection{响应式断点}

系统面向桌面端使用，但要求在不同屏幕尺寸下保持可用：

\begin{itemize}
    \item \textbf{大屏（$\ge$1200px）}：三栏布局完整展开，适合在1920×1080显示器上使用；
    \item \textbf{中屏（768-1199px）}：右侧面板自动收起为浮动按钮，点击后以Drawer形式展开；
    \item \textbf{小屏（<768px）}：切换为单栏布局，侧边导航自动收起为汉堡菜单。
\end{itemize}

\subsection{组件使用约定}

\begin{itemize}
    \item \textbf{表格（el-table）}：数据量超过20条时分页，每页20条。表格行支持hover高亮，关键操作的按钮在操作列展示；
    \item \textbf{表单（el-form）}：必填字段标红色星号，提交前做前端校验，校验不通过时在字段下方显示错误提示文字；
    \item \textbf{对话框（el-dialog）}：宽度统一为600px或800px，关闭方式支持点击遮罩关闭和右上角X关闭。提交操作后自动关闭并刷新父页面数据；
    \item \textbf{消息提示（ElMessage）}：操作成功使用绿色提示（成功），操作失败使用红色提示（错误），持续时间均为3秒；
    \item \textbf{加载状态（v-loading）}：数据获取和表单提交过程中显示加载遮罩，防止重复操作。
\end{itemize}
